\chapter{Literature Review}
\label{ch:lit_rev}

This chapter reviews the theoretical and algorithmic foundations of reinforcement 
learning as studied in the course material provided by Prof.\ Hirchoua Badr 
at ENSAM Casablanca. It situates this project within the current state of the 
art, identifies the gap that motivates the chosen approach, and demonstrates 
the significance of applying multi-agent deep reinforcement learning to a 
physically realistic football environment. The review is organised thematically: 
Section~\ref{sec:lit_rl_foundations} covers the core RL formalism; 
Section~\ref{sec:lit_value_based} reviews value-based methods; 
Section~\ref{sec:lit_policy_based} reviews policy-based and actor-critic methods; 
Section~\ref{sec:lit_marl} addresses multi-agent RL; 
Section~\ref{sec:lit_critique} identifies the gap; and 
Section~\ref{sec:lit_summary} summarises the chapter.

% ================================================================
\section{Reinforcement Learning Foundations}
\label{sec:lit_rl_foundations}
% ================================================================

Reinforcement learning is a computational framework for sequential 
decision-making under uncertainty~\citep{sutton2018reinforcement}. Unlike 
supervised learning, which requires labelled input-output pairs, RL relies 
solely on a reward signal obtained through interaction. The agent observes a 
state $s_t$, selects an action $a_t$ according to its policy $\pi$, and receives 
a scalar reward $r_t$ from the environment, transitioning to a new state $s_{t+1}$.

\subsection{Markov Decision Process}

The standard mathematical framework for RL is the Markov Decision Process (MDP), 
defined by the tuple $(\mathcal{S}, \mathcal{A}, \mathcal{P}, \mathcal{R}, \gamma)$, 
where $\mathcal{S}$ is the state space, $\mathcal{A}$ is the action space, 
$\mathcal{P}(s'|s,a)$ is the state transition probability, $\mathcal{R}(s,a)$ 
is the expected reward, and $\gamma \in [0,1)$ is the discount factor. 
The Markov property requires that the next state depends only on the current 
state and action, not on prior history:
\begin{equation}
    \mathbb{P}[S_{t+1} \mid S_t] = \mathbb{P}[S_{t+1} \mid S_1, \ldots, S_t].
    \label{eq:markov}
\end{equation}

\subsection{Return and Value Functions}

The goal of the agent is to maximise the expected discounted cumulative reward, 
known as the return:
\begin{equation}
    G_t = \sum_{k=0}^{\infty} \gamma^k r_{t+k+1}.
    \label{eq:return}
\end{equation}

The state-value function $V^\pi(s)$ and the action-value function $Q^\pi(s,a)$ 
measure how good it is to be in a given state or to take a given action under 
policy $\pi$:
\begin{equation}
    V^\pi(s) = \mathbb{E}_\pi \left[ G_t \mid S_t = s \right], \quad
    Q^\pi(s,a) = \mathbb{E}_\pi \left[ G_t \mid S_t = s, A_t = a \right].
    \label{eq:value_functions}
\end{equation}

\subsection{The Bellman Equation}

The Bellman equation expresses the recursive relationship at the heart of RL. 
For the optimal value function~\citep{bellman1957dynamic}:
\begin{equation}
    V^*(s) = \max_{a} \left[ \mathcal{R}(s,a) + \gamma \sum_{s'} 
    \mathcal{P}(s'|s,a) V^*(s') \right].
    \label{eq:bellman}
\end{equation}
This equation underpins all dynamic programming and temporal-difference methods. 
In the tabular Q-Learning demo (Section~\ref{sec:results_qlearning}), the agent 
directly applies this principle through iterative table updates.

% ================================================================
\section{Value-Based Methods}
\label{sec:lit_value_based}
% ================================================================

Value-based methods learn an estimate of the action-value function $Q(s,a)$ 
and derive a policy implicitly by acting greedily with respect to it.

\subsection{Q-Learning}

Q-Learning~\citep{watkins1992qlearning} is an off-policy temporal-difference 
algorithm that updates the Q-table using the Bellman optimality equation:
\begin{equation}
    Q(s,a) \leftarrow Q(s,a) + \alpha 
    \left[ r + \gamma \max_{a'} Q(s', a') - Q(s,a) \right],
    \label{eq:qlearning}
\end{equation}
where $\alpha$ is the learning rate. Action selection follows the 
$\varepsilon$-greedy policy: with probability $\varepsilon$ the agent 
explores randomly, and with probability $1 - \varepsilon$ it exploits 
the best known action. Q-Learning is guaranteed to converge to the optimal 
policy under mild conditions~\citep{watkins1992qlearning}. This algorithm 
was implemented in the \texttt{q\_learning\_demo.py} script on FrozenLake-v1.

\subsection{Deep Q-Network (DQN)}

The key limitation of tabular Q-Learning is its inability to scale to large 
or continuous state spaces. \cite{mnih2015humanlevel} addressed this by 
replacing the Q-table with a deep neural network $Q(s,a;\theta)$, parameterised 
by weights $\theta$. Two critical innovations were introduced:

\begin{itemize}
    \item \textbf{Experience Replay:} Transitions $(s, a, r, s')$ are stored in 
    a replay buffer $\mathcal{D}$, and mini-batches are sampled uniformly at 
    random for training, breaking temporal correlations between consecutive samples.
    
    \item \textbf{Target Network:} A separate network $Q(s,a;\theta^-)$ with 
    frozen weights is used to compute TD targets, improving training stability:
    \begin{equation}
        y_i = r + \gamma \max_{a'} Q(s', a'; \theta^-).
        \label{eq:dqn_target}
    \end{equation}
\end{itemize}

The DQN loss is the mean squared error between the online network's prediction 
and the target:
\begin{equation}
    \mathcal{L}(\theta) = \mathbb{E}_{(s,a,r,s') \sim \mathcal{D}} 
    \left[ \left( y_i - Q(s,a;\theta) \right)^2 \right].
    \label{eq:dqn_loss}
\end{equation}

This algorithm was implemented in \texttt{dqn\_demo.py} on CartPole-v1, 
as covered in Lab 2 of the course.

% ================================================================
\section{Policy-Based and Actor-Critic Methods}
\label{sec:lit_policy_based}
% ================================================================

Value-based methods are limited to discrete action spaces and cannot 
represent stochastic policies. Policy-based methods directly parameterise 
the policy $\pi_\theta(a|s)$ and optimise it with respect to the 
expected return~\citep{sutton1999policy}:
\begin{equation}
    J(\pi_\theta) = \mathbb{E}_{\tau \sim \pi_\theta} \left[ R(\tau) \right],
    \label{eq:policy_objective}
\end{equation}
where $\tau = (s_0, a_0, r_1, \ldots)$ is a trajectory.

\subsection{Policy Gradient Theorem and REINFORCE}

The policy gradient theorem~\citep{sutton1999policy} provides an analytical 
expression for the gradient of $J$ without requiring knowledge of the 
environment dynamics:
\begin{equation}
    \nabla_\theta J(\pi_\theta) = \mathbb{E}_{\tau \sim \pi_\theta} 
    \left[ \sum_{t=0}^{T} \nabla_\theta \log \pi_\theta(a_t|s_t) \cdot R(\tau) \right].
    \label{eq:policy_gradient}
\end{equation}
The REINFORCE algorithm~\citep{williams1992reinforce} implements this estimator 
using Monte Carlo rollouts. Its main drawback is high variance, addressed by 
actor-critic methods.

\subsection{Proximal Policy Optimisation (PPO)}

PPO~\citep{schulman2017proximal} is currently one of the most widely used 
deep RL algorithms for continuous control. It improves upon earlier methods 
by constraining policy updates to a trust region, preventing destructively 
large gradient steps. The PPO clipped objective is:
\begin{equation}
    \mathcal{L}^{\text{CLIP}}(\theta) = \mathbb{E}_t \left[ 
    \min\left( r_t(\theta) \hat{A}_t,\; 
    \text{clip}(r_t(\theta), 1-\varepsilon, 1+\varepsilon) \hat{A}_t 
    \right) \right],
    \label{eq:ppo_clip}
\end{equation}
where $r_t(\theta) = \frac{\pi_\theta(a_t|s_t)}{\pi_{\theta_\text{old}}(a_t|s_t)}$ 
is the probability ratio, $\hat{A}_t$ is the estimated advantage, and 
$\varepsilon$ is the clipping coefficient. PPO is the algorithm used in 
the main contribution of this project.

% ================================================================
\section{Multi-Agent Reinforcement Learning}
\label{sec:lit_marl}
% ================================================================

In multi-agent settings, multiple agents learn simultaneously within a shared 
environment. This introduces non-stationarity: from any single agent's perspective, 
the environment appears non-Markovian because other agents' policies change 
during training~\citep{lowe2017multiagent}.

\subsection{Parameter Sharing}

Parameter sharing is an approach where all agents share a single policy 
network~\citep{gupta2017cooperative}. Each agent observes its own local state 
and acts independently at execution time, but all agents contribute to the same 
gradient update during training. This substantially reduces the number of 
parameters to learn and has been shown to be effective in symmetric multi-agent 
tasks. In this project, parameter sharing is implemented via SuperSuit's 
\texttt{concat\_vec\_envs} utility, which treats all agents' observations 
as independent parallel environments feeding into a single PPO policy.

\subsection{Centralised Training, Decentralised Execution (CTDE)}

The CTDE paradigm~\citep{lowe2017multiagent} allows agents to access global 
information during training (centralised critic) while acting on local 
observations during execution (decentralised actors). While MADDPG and MAPPO 
implement full CTDE, this project uses the simpler parameter-sharing variant 
as a practical first step toward multi-agent coordination.

% ================================================================
\section{Critique of the Review}
\label{sec:lit_critique}
% ================================================================

The reviewed methods form a clear progression from tabular methods 
(Q-Learning) to function approximation (DQN) and finally to policy 
optimisation (PPO). The key gap addressed by this project is the 
application of these techniques to a continuous, physically realistic, 
multi-agent environment. Most course examples are restricted to discrete, 
grid-world environments (FrozenLake, CartPole). Applying PPO to MuJoCo Soccer 2v2 
represents a significant increase in complexity: the state and action spaces 
are continuous, the environment involves coupled multi-agent dynamics, and 
reward signals are naturally sparse. The reward shaping strategy developed 
in this project (Section~\ref{sec:rw_shaping}) directly addresses the 
sparse reward gap, a limitation not covered in the existing literature 
on these specific environments.

% ================================================================
\section{Summary}
\label{sec:lit_summary}
% ================================================================

This chapter reviewed the existing literature relevant to the project, 
covering the MDP formalism, Bellman equations, Q-Learning, DQN, policy 
gradients, PPO, and multi-agent RL. The review identified that state-of-the-art 
algorithms such as PPO, combined with parameter sharing and reward shaping, 
provide the most appropriate framework for the cooperative and competitive 
multi-agent football problem investigated here. Chapter~\ref{ch:method} 
describes the methodology used to apply these algorithms to the MuJoCo 
Soccer 2v2 environment.